% Chapter 2, section 9: Ephemeral queues
% Write in plain LaTeX. Code of a guide step: \guide{<step>}; commands: \begin{shell}...\end{shell};
% a screenshot: \screenshot{file}{caption}; a gap in the draft: \placeholder{...}.
% See book/latex/queen.sty for the full list.

\section{Ephemeral queues}

Ephemeral queues live in memory, outside the log, for request/reply. The worker answers on whatever
inbox a request names; the caller makes a private inbox and waits on it. There is no transaction or
replay here, and a crash loses what is in flight.

The worker:

\guide{worker}

The caller:

\guide{caller}

\why{a request that waits for an answer (``what is the price of A-1, now?'') needs speed more
than history. Writing it to disk, replicating it and keeping it for days would cost time for
nothing. Ephemeral queues keep the same groups, leases and acks, and skip the log. Use a regular
queue for anything that must survive a crash.}

\queenfig{ephemeral-reply}{Request and reply over ephemeral queues, with a private inbox.}

On the site, with the code in all six SDK languages: \href{https://queenmq.com/writing-applications/ephemeral/}{queenmq.com/writing-applications/ephemeral}.
